\chapter{Realism, Anti-Realism, and the Problem of Success}
\label{ch:realism}

\begin{chapterguide}
This chapter asks what the success of science licenses us to believe. It
examines constructive empiricism, scientific realism, structural realism,
entity realism, pragmatism, inference to the best explanation, and the
pessimistic meta-induction. The chapter's conclusion is a layered conception of
truth and a graded account of ontological commitment.
\end{chapterguide}

The central realist argument is simple. Mature scientific theories often
predict novel phenomena, integrate diverse observations, guide new instruments,
and support interventions. It would be a remarkable coincidence if the
theoretical entities and structures responsible for this success were entirely
fictional. Scientific realism therefore recommends belief in at least some
unobservable aspects of the world described by our best theories
\citep{chakravartty2007}.

The anti-realist replies that success need not reveal truth. A theory can be
empirically adequate: it can correctly describe what is observable while
remaining agnostic about unobservable entities. Bas van Fraassen's constructive
empiricism is not a claim that electrons do not exist. It is a claim about the
epistemic attitude required by science: acceptance may involve belief that a
theory is empirically adequate without belief that its unobservable posits are
true \citep{vanfraassen1980}.

Both positions contain an important distinction. The question is not whether a
scientist may use a model containing unobservables. Scientists plainly do.
The question is what attitude the use and success of the model rationally
warrant.

\section{Constructive empiricism}

Constructive empiricism avoids a crude instrumentalism. It can accept that
scientific theories are true-or-false claims, that they explain phenomena, and
that scientists can rationally adopt them for practical purposes. Its
distinctive requirement is epistemic modesty: belief need extend only to
observable adequacy.

This view has several strengths. It respects the historical fact that
successful theories are often replaced. It avoids treating every useful
mathematical posit as an entity. It separates the aim of science from the
psychology of individual scientists. And it recognizes that seeing through a
microscope or detector involves a complex notion of observability rather than
an absolute line between the visible and invisible.

Its difficulty is that the observable/unobservable distinction can be unstable
and method-dependent. A virus, a distant galaxy, and a magnetic field do not
form a natural epistemic class. More importantly, empirical adequacy can be
achieved by models whose internal structure supports counterfactual and
interventional reasoning. If a theory's unobservable components are repeatedly
used to make new kinds of contact with the world, permanent agnosticism may be
too restrictive.

\section{Scientific realism and the no-miracles intuition}

Realists often appeal to the no-miracles argument: the ability of science to
predict and intervene would be extraordinary if theories were not at least
approximately true. The argument is not a deductive proof. It is an inference
to the best explanation (IBE)\index{inference to the best explanation}. We
compare explanations of scientific success and regard truth-tracking as one
candidate.

Gilbert Harman used the phrase \enquote{inference to the best explanation} for
the pattern in which a hypothesis is selected because it would, if true,
explain the evidence better than alternatives \citep{harman1965}. IBE is not a
license to accept whatever story feels satisfying. A good explanatory
hypothesis should have independently supported consequences, expose itself to
failure, and outperform rivals in a specified domain.

The no-miracles argument is strongest when success is robust across
independent routes: a theory predicts an effect before it is observed, survives
instrumental changes, connects different phenomena, and supports interventions
that would not be possible if the relevant structure were absent. It is weaker
when the success is merely a fit to existing data, depends on flexible
parameters, or is achieved by many unrelated models.

\section{The pessimistic meta-induction}

The historical record supplies an obvious challenge. Many past theories were
successful and yet contained entities or principles that later scientists
rejected. The luminiferous ether supported a useful wave theory of light.
Phlogiston organized chemical phenomena before oxygen chemistry. Classical
mechanics remains reliable in many domains although its ontology is not the
final physical description.

The pessimistic meta-induction argues that our current theories may be like
those past theories: successful for a time but false in their unobservable
commitments. This is not a proof of anti-realism. It is a warning against
treating present success as a guarantee of total truth.

The realist response often distinguishes between abandoned structures and
retained structures. A later theory may reject an entity while preserving
equations, symmetry relations, limiting cases, or causal dependencies. John
Worrall's structural realism emphasizes precisely this possibility
\citep{worrall1989}. The theory changes, but some relations survive.

The response is plausible but incomplete. Mathematical structure is not a
single object. One must specify which structure is retained, at what level, and
under what transformations. A merely formal similarity can survive while the
physical interpretation changes. \RCR\ therefore treats structural continuity
as evidence, not as an automatic guarantee of truth.

\section{Structural realism}

Structural realism comes in epistemic and ontic forms. Epistemic structural
realism says that science gives us knowledge of relations or structure even if
knowledge of the intrinsic nature of entities remains limited. Ontic structural
realism makes the stronger metaphysical claim that structure is more
fundamental than objects.

The epistemic version fits an important pattern in theory change. A relation
such as a mathematical dependency may survive when the entities used to
interpret it change. The relational pattern can therefore be a better target of
realist commitment than the full ontology of a theory.

The ontic version is more ambitious but faces questions about what a structure
is without relata, how structures are individuated, and how causal powers fit
into a purely relational picture. \RCR\ remains neutral between object-first
and structure-first metaphysics at the ultimate level. It claims only that
scientific evidence often supports structured constraints more securely than
an exhaustive inventory of intrinsic natures.

\section{Entity realism and intervention}

Ian Hacking shifted attention from representation to intervention. When
scientists can use an entity to create stable effects in other systems, there
is a strong practical reason to believe that the entity is not merely a
calculational fiction \citep{hacking1983}. This is especially persuasive in
experimental settings where the entity is not just inferred from the same data
that the theory was designed to fit.

Intervention is not infallible. A successful manipulation can be described by
different ontologies. An instrument may be responding to a field, a mechanism,
or a composite process. Experimental control can also be achieved through a
model whose components are only approximately isolated.

The lesson is therefore conditional. Intervention raises the ontological
status of a posit when the effect is robust under changes in apparatus,
calibration, and background assumptions, and when the posit makes a difference
to outcomes that were not used to define it. This will become one of the
criteria in \RCR's commitment ladder.

\section{Pragmatism without relativism}

Pragmatist traditions, especially Peirce's account of inquiry, emphasize that
belief is a habit of action and that the meaning of an idea lies partly in its
conceivable practical consequences \citep{peirce1878}. This is not the claim
that truth means whatever works for a person. Peirce's ideal of inquiry is
public, fallible, and oriented toward a community that could continue testing
the claim.

Pragmatism corrects two realist errors. First, it reminds us that theories are
tools for asking questions and changing situations, not museum labels attached
to a finished world. Second, it prevents the search for truth from becoming
detached from the practices that give concepts their use.

It also needs correction. Immediate usefulness can reward a false model,
especially when a short-term prediction is obtained through a fragile
correlation. \RCR\ distinguishes pragmatic adequacy from correspondence and
causal adequacy. A model can be useful while its ontological interpretation
remains uncertain.

\section{A layered account of truth}

The debate becomes less confusing when different kinds of scientific success
are separated.

\begin{table}[ht]
\centering
\smalltable
\begin{tabularx}{\textwidth}{p{0.22\textwidth}Y Y}
\toprule
Layer & Question & Typical evidence\\
\midrule
Correspondence & Does the claim identify what exists or occurs? & Independent detection, historical trace, convergent measurement\\
Structural accuracy & Does the representation preserve relevant relations? & Invariance, limiting cases, cross-domain integration\\
Interventional adequacy & Does changing the represented variable change outcomes as expected? & Experiment, manipulation, mechanism, counterfactual test\\
Pragmatic adequacy & Does the representation support reliable inquiry or action? & Forecasting, control, diagnosis, design\\
\bottomrule
\end{tabularx}
\caption{Four layers of scientific success. They support one another but do not
collapse into one another.}
\label{tab:truth-layers}
\end{table}

A model may be structurally accurate in a restricted regime while being
literally false outside it. A variable may be useful for intervention without
being a fundamental entity. A theory may correspond to a real mechanism while
its mathematical idealization misstates the mechanism's fine detail. The
proper question is not whether the whole theory is simply true, but which
claims, structures, and interventions have earned which kinds of commitment.

\begin{principlebox}
\textbf{Layered truth principle.} Scientific truth is not one undifferentiated
property. A claim may be correct about a phenomenon, structurally correct
about a relation, causally adequate for an intervention, and pragmatically
useful to different degrees. Objectivity improves when these layers are
reported separately rather than compressed into one label.
\end{principlebox}

\section{What realism should commit us to}

\RCR\ rejects both the demand for literal truth of every model component and
the refusal to believe in any unobservable. It recommends graded commitment.

\begin{enumerate}
\item Believe an observable event when its measurement path is reliable and
alternative explanations have been tested.
\item Treat a structure as objectively supported when it remains stable across
independent representations and transformations.
\item Treat an entity or mechanism as strongly supported when intervention,
convergence, and counterfactual robustness all point to it.
\item Treat a theoretical ontology as provisional when its success is confined
to prediction or depends on highly flexible idealization.
\end{enumerate}

This is not an attempt to turn realism into a numerical score. It is a
discipline of separating what the evidence supports from what the imagination
adds. The final commitment should be proportional to the strongest constraint
that the inquiry has actually established.

\section*{Chapter summary}

Constructive empiricism protects epistemic modesty; scientific realism explains
why robust success can support belief in unobservables; structural realism
highlights continuity of relations; entity realism emphasizes intervention; and
pragmatism links meaning to public inquiry and action. The pessimistic
meta-induction warns that successful theories can contain false commitments.
The proposed response is layered truth and graded commitment: correspondence,
structural, interventional, and pragmatic adequacy are distinct dimensions of
scientific success.

\begin{discussion}
\begin{enumerate}
\item Is an entity more credible when it can be manipulated, or can
intervention itself be theory-dependent?
\item Which parts of a modern scientific theory seem most secure: entities,
equations, structures, mechanisms, or measurement procedures?
\item Can empirical adequacy be enough for practical science even when
ontology is unsettled?
\end{enumerate}
\end{discussion}

\begin{furtherreading}
Read van Fraassen (1980), Hacking (1983), Worrall (1989), Fine (1984),
Chakravartty (2007), and Stanford (2006) as contrasting positions in the
realism debate.
\end{furtherreading}

